\chapter{Родитель коэффициента связывания композитного
\texorpdfstring{$SU(2)$}{SU(2)}-синглета}
\label{ch:v10-discrete-resolution-su2-singlet-binding-coefficient}

Предыдущий аудит выделил канонический обменный оператор
\begin{equation}
 X=\sum_aT_a\otimes T_a=\frac14I-P_S,
\end{equation}
но не вывел коэффициент перед ним. Унаследованное значение $g^2=3/8$
является коэффициентом вершины, а не готовой энергией связывания: для
перехода к двухчастичному gap требуется статическая восприимчивость
посредника $\chi$.

Минимальный условный родитель имеет вид
\begin{equation}
 \mathcal P(\kappa,\chi)=\frac12(\kappa-g^2\chi)^2
 +\frac12(\chi-\chi_0)^2.
\end{equation}
Его стационарная точка равна
\begin{equation}
 (\kappa_*,\chi_*)=(g^2\chi_0,\chi_0)
 =\left(\frac38\chi_0,\chi_0\right),
\end{equation}
а гессиан
\begin{equation}
 H=\begin{pmatrix}1&-3/8\\-3/8&73/64\end{pmatrix},
 \qquad \det H=1,
\end{equation}
строго положителен. Условный оператор $\kappa_*P_T$ оставляет синглет при
нуле и поднимает триплет на $\kappa_*$.

Однако унаследованный родитель на координатах $(\kappa,\chi)$ имеет
гессиан $\operatorname{diag}(1,0)$: он фиксирует только
$\kappa=0$, оставляет $\chi$ плоской и не содержит смешанного блока
$-g^2$. Поэтому число $3/8$ не превращается в binding gap без нового
пропагатора или восприимчивости. Кроме того, отсутствует отображение от
$\kappa_*$ к массе композитного полюса.

\begin{theorem}[Условный коэффициент имеет строгий родитель]
Если независимо задана восприимчивость $\chi_0$, положительный общий
родитель однозначно выводит $\kappa_*=g^2\chi_0$ и синглет-триплетное
расщепление. Архитектура совместима с каноническим обменом.
\end{theorem}

\begin{nogocriterion}[Вершина не равна связанному состоянию]
Значение $g^2=3/8$ не является коэффициентом потенциала без типизированного
посредника, его пропагатора и выбранного статического предела. Подстановка
$\chi_0=1$ была бы выбором единиц, а не физическим происхождением gap.
\end{nogocriterion}

\begin{protocol}\begin{enumerate}
\item условная архитектура: \textbf{$10/10$};
\item унаследованные кинематические ингредиенты: \textbf{$2/2$};
\item унаследованный динамический mixed block: \textbf{$0/1$};
\item унаследованный родитель: \textbf{ранг/ядро $1/1$};
\item физическое происхождение: \textbf{$0/3$};
\item следующий гейт: \textbf{аудит восприимчивости посредника}.
\end{enumerate}\end{protocol}

Машинный сертификат сохранён в
\path{s2t_v10_cell_birth_four_volume_nonequilibrium_bath_discrete_resolution_transmuted_mode_asymptotically_free_su2_gauge_singlet_binding_coefficient_parent_origin_gate_results.json}.